\documentclass[11pt]{article}
\usepackage[utf8]{inputenc}
\usepackage[T1]{fontenc}
\usepackage{lmodern}
\usepackage[spanish]{babel}
\usepackage[margin=2.5cm]{geometry}
\usepackage{url}

\title{Apertura a la descentralización:\\ más allá de los pesos abiertos}
\author{Laura Itzel Rodríguez Dimayuga}
\date{Noviembre 2026}

\begin{document}

\maketitle

\section*{Tesis central}

La apertura de un gran modelo de lenguaje (LLM) no consiste en poder descargar sus pesos. Consiste en que las comunidades puedan usar, estudiar, modificar y auditar los sistemas que median su conocimiento, sus decisiones y sus relaciones. De ahí se sigue una tesis más fuerte: la etiqueta «abierto», aplicada a los modelos de peso abierto, describe una forma de distribución y no una redistribución de poder. Quien recibe los pesos puede ejecutar el modelo, pero no puede reconstruir cómo se hizo ni con qué datos, y tampoco puede producir otro. La descentralización que promete la apertura solo ocurre cuando se abren también los datos de entrenamiento, el código y el proceso, y cuando existen condiciones materiales para aprovecharlos. Por eso la disputa sobre qué cuenta como «abierto» es política antes que técnica: decide quién controla los medios de producción cognitivos.

\section*{Problema}

Sapkota et al.\ (2025) examinan más de cien modelos recientes y encuentran que muchos de los que se anuncian como de código abierto no reportan sus datos ni su código de entrenamiento, práctica que la literatura llama \textit{open-washing}. OLMo (Groeneveld et al., 2024) y Pythia (Biderman et al., 2023) muestran que la apertura completa es practicable, porque publican datos, código de entrenamiento y puntos de control intermedios .

La regulación reproduce la ambigüedad. El Reglamento de IA de la Unión Europea exime de ciertas obligaciones de transparencia a los modelos de propósito general con licencia libre cuyos pesos, arquitectura e información de uso sean públicos. El mismo texto admite que esa liberación no revela información sustancial sobre los datos de entrenamiento, y por ello mantiene la obligación de publicar un resumen del contenido usado y una política de derechos de autor. Qiu et al.\ (2026) parten de esa indefinición y modelan cómo el umbral de apertura que fija el regulador altera las decisiones de publicación de los desarrolladores.

\section*{Estrategia argumentativa}

El argumento avanza en cinco pasos.

\begin{enumerate}
\item \textbf{Distinción conceptual.} Separo peso abierto de apertura completa y propongo cuatro capacidades como criterio: usar, estudiar, modificar y auditar. Los pesos garantizan la primera y solo parcialmente las otras tres.

\item \textbf{Lectura crítica de la regulación.} Muestro que el Reglamento europeo define la apertura por los artefactos publicados y no por las capacidades que habilita, y que reconoce esa limitación al conservar las obligaciones sobre datos.

\item \textbf{Criterio normativo.} Manchanda et al.\ (2026) sostienen, desde el enfoque de capacidades, que el desarrollo abierto amplía la transparencia, la reproducibilidad, la participación comunitaria y la auditabilidad. Nasi (2026) propone que una IA es democrática cuando redistribuye poder, tanto en su elaboración como en su acción. Con ambos argumento que el peso abierto interviene en el uso del modelo y deja intacta su elaboración.

\item \textbf{Situación.} Leo el problema desde América Latina con la teoría marxista de la dependencia (Reyes Núñez, 2025).

\item \textbf{Objeciones.} Respondo a tres: que abrir datos y procesos aumenta el riesgo de uso indebido; que los modelos completamente abiertos aún rinden menos, como reconocen Manchanda et al.; y que abrir los datos no descentraliza nada si el cómputo sigue concentrado. Concedo la tercera en parte: la apertura completa es condición necesaria de la descentralización, no suficiente.
\end{enumerate}

\section*{Posición respecto a los ejes}

El artículo se sitúa principalmente en los ejes de infraestructura y conocimiento, y recurre a los de control y relaciones como evidencia de lo que está en juego.

\paragraph{Infraestructura (Marx).} Datos, cómputo, código y modelos son medios de producción. Reyes Núñez (2025) describe la dependencia tecnológica como rasgo estructural del capitalismo latinoamericano, con participación marginal en la producción de tecnologías de vanguardia y transferencias de valor hacia el exterior. Un modelo de peso abierto encaja en ese patrón: la región lo consume y lo ajusta, pero no adquiere la capacidad de producirlo. El número 559 de la revista de ALAI (2026) aporta la discusión regional sobre infraestructuras y soberanía digital.

\paragraph{Conocimiento (Habermas).} Manchanda et al.\ (2026) tratan a los LLMs como infraestructura que media el acceso al conocimiento, y Machado (2025) asocia los modelos abiertos con escrutinio independiente y mitigación de sesgos. Si la deliberación pública pasa por estos sistemas, no puede deliberarse sobre lo que no puede inspeccionarse.

\paragraph{Control (Foucault).} Nasi (2026) retoma a Foucault para pensar el poder como relación presente en la vida cotidiana. De Jager y Albarracin (2026) lo describen como capacidad de modular la atención mediante guiones compartidos. Holmov et al.\ (2026) muestran que los LLMs infieren atributos de sus usuarios y condicionan su conducta a esas inferencias, incluidas las negativas a responder, y que ese estado interno puede leerse y modificarse. Este caso matiza la tesis: esa auditoría ya es posible con acceso a los pesos. Lo que los pesos no dan es la procedencia de los datos ni la posibilidad de rehacer el modelo.

\paragraph{Relaciones.} Serapio-García et al.\ (2023) muestran que la personalidad que exhibe un LLM puede medirse y moldearse. Maeda y Quan-Haase (2024) analizan cómo el diseño antropomórfico induce confianza y vínculos parasociales. Quien controla el modelo controla también el tono del vínculo, y los usuarios de un sistema cerrado no pueden verificar con qué fines se ajusta.

\section*{Contribución esperada}

El artículo ofrece un criterio de apertura basado en capacidades comunitarias y una lectura situada, desde la dependencia tecnológica, de por qué el peso abierto no basta.

\section*{Referencias preliminares}

\begingroup
\setlength{\parindent}{-1.5em}
\setlength{\leftskip}{1.5em}
\setlength{\parskip}{0.4em}
\raggedright

ALAI (2026). \textit{Pax Silica: nueva estrategia de guerra}, núm.\ 559, febrero. Incluye C.~Reyes, «Las infraestructuras de la IA: poder corporativo y polarización global», y M.~E.~Stédile, «Soberanía digital en el Sur Global».

Biderman, S. et al.\ (2023). Pythia: A suite for analyzing large language models across training and scaling. \textit{International Conference on Machine Learning (ICML)}, 2397--2430.

de Jager, S. y Albarracin, M. (2026). Post-control script-societies: Social power dynamics through attention mechanisms as active inference. \textit{Topoi}. doi:10.1007/s11245-026-10458-7

Groeneveld, D. et al.\ (2024). OLMo: Accelerating the science of language models. \textit{Proceedings of ACL 2024}, 15789--15809. doi:10.18653/v1/2024.acl-long.841

Holmov, A., Huang, Y., Bykov, K. y Akata, Z. (2026). User model extraction via belief self-distillation. arXiv:2609.31603.

Machado, J. (2025). Toward a public and secure generative AI: A comparative analysis of open and closed LLMs. arXiv:2505.10603.

Maeda, T. y Quan-Haase, A. (2024). When human-AI interactions become parasocial: Agency and anthropomorphism in affective design. \textit{FAccT '24}.

Manchanda, J., Boettcher, L., Westphalen, M. y Jasser, J. (2026). The open-source advantage in large language models (LLMs). \textit{AI and Ethics}, 6, 390.

Nasi, F. (2026). When is AI democratic? Artificial intelligence and democratic empowerment. \textit{AI \& Society}, 41, 5481--5494.

Parlamento Europeo y Consejo de la Unión Europea (2024). Reglamento (UE) 2024/1689 (Reglamento de Inteligencia Artificial). \textit{Diario Oficial de la Unión Europea}, serie L.

Qiu, T., Laufer, B., Kleinberg, J. y Heidari, H. (2026). Modeling the economic impacts of AI openness regulation. \textit{Advances in Neural Information Processing Systems}.

Reyes Núñez, C. (2025). Patrón de reproducción del capital y dependencia tecnológica. \textit{Estudios Latinoamericanos}, 39(54), 141--163. doi:10.22201/fcpys.24484946e.2024.54.93745

Sapkota, R., Raza, S. y Karkee, M. (2025). Comprehensive analysis of transparency and accessibility of ChatGPT, DeepSeek, and other SoTA large language models. arXiv:2502.18505.

Serapio-García, G. et al.\ (2023). Personality traits in large language models. arXiv:2307.00184.

\endgroup

\end{document}
